\section*{छोटका रूप आ संकेत --- Abbreviations \& Symbols}

\vspace{0.2em}
{\small\color{bpgray}Grammatical, linguistic origin labels, and frequency indicators used throughout this trilingual Bhojpuri dictionary.}
\vspace{0.6em}

\begin{multicols}{2}
\raggedcolumns

% --- Section 1: Word Classes ---
{\color{bpblue}\bfseries\small शब्द के जात / Word Classes}
\vspace{0.2em}

\begin{description}[leftmargin=1.0em, labelwidth=!, itemsep=0.1em, font=\bfseries\color{bpbluedark}]
  \item[संज्ञा (n.)] Noun / नाम (नाँव)
  \item[सर्वनाम (pron.)] Pronoun / नाम के बदला
  \item[विशेषण (adj.)] Adjective / गुन-लछन
  \item[क्रिया (v.)] Verb / काम-काज
  \item[क्रिया विशेषण (adv.)] Adverb
  \item[परसर्ग (postp.)] Postposition / पाछा लागल शब्द
  \item[संयोजक (conj.)] Conjunction / जोड़ेवाला
  \item[अव्यय (part.)] Particle / अटल शब्द
  \item[विस्मयादिबोधक (interj.)] Interjection / अचरज बोधक
\end{description}

\vspace{0.5em}

% --- Section 2: Origin & Frequency ---
{\color{bpblue}\bfseries\small शब्द के स्रोत आ प्रयोग / Origin \& Frequency}
\vspace{0.2em}

\begin{description}[leftmargin=1.0em, labelwidth=!, itemsep=0.1em, font=\bfseries\color{bpbluedark}]
  \item[ठेठ भोजपुरी] Pure Native Regional Spoken Terms
  \item[तद्भव] Vernacular Words Evolved from Sanskrit
  \item[तत्सम/साहित्यिक] Classical Sanskrit Literary Loanwords
  \item[फारसी/अरबी आगत] Perso-Arabic Loanwords in Bhojpuri
  \item[\freqcommon] Common Everyday Vocabulary
  \item[\freqformal] Formal or Literary Vocabulary
  \item[\freqrare] Rare, Archaic, or Cultural Terms
\end{description}

\columnbreak

% --- Section 3: Grammatical Features ---
{\color{bpblue}\bfseries\small व्याकरणिक भेद / Features}
\vspace{0.2em}

\begin{description}[leftmargin=1.0em, labelwidth=!, itemsep=0.1em, font=\bfseries\color{bpbluedark}]
  \item[लिंग (g.)] Gender / जात
  \item[पुल्लिंग (m.)] Masculine / मरदाना
  \item[स्त्रीलिंग (f.)] Feminine / जनाना
  \item[वचन (num.)] Number / गिनती
  \item[एकवचन (sg.)] Singular / एगो
  \item[बहुवचन (pl.)] Plural / अनेक
  \item[पुरुष (pers.)] Person / पुरख
  \item[पहिल पुरख (1p.)] 1st Person / उत्तम पुरुष
  \item[दूसर पुरख (2p.)] 2nd Person / मध्यम पुरुष
  \item[तीसर पुरख (3p.)] 3rd Person / अन्य पुरुष
\end{description}

\vspace{0.5em}

% --- Section 4: TAM (Tense-Aspect-Mood) ---
{\color{bpblue}\bfseries\small काल आ भाव / Tense \& Aspect}
\vspace{0.2em}

\begin{description}[leftmargin=1.0em, labelwidth=!, itemsep=0.1em, font=\bfseries\color{bpbluedark}]
  \item[चलते काल (pres.)] Present Tense / वर्तमान
  \item[बीतल काल (past)] Past Tense / भूत
  \item[आवेवाला काल (fut.)] Future Tense / भविष्य
  \item[पूरन काल (perf.)] Perfective Aspect / पूर्ण
  \item[अधुरा काल (imperf.)] Imperfective Aspect / अपूर्ण
  \item[हुकुम/आज्ञा (imp.)] Imperative Mood
  \item[इच्छा/चाहत (opt.)] Optative / Subjunctive
\end{description}

\end{multicols}
